\documentclass[12pt, titlepage]{article}

\usepackage{amsmath, mathtools}

\usepackage[round]{natbib}
\usepackage{amsfonts}
\usepackage{amssymb}
\usepackage{graphicx}
\usepackage{colortbl}
\usepackage{xr}
\usepackage{hyperref}
\usepackage{longtable}
\usepackage{xfrac}
\usepackage{tabularx}
\usepackage{float}
\usepackage{siunitx}
\usepackage{booktabs}
\usepackage{multirow}
\usepackage[section]{placeins}
\usepackage{caption}
\usepackage{fullpage}

\hypersetup{
bookmarks=true,     % show bookmarks bar?
colorlinks=true,       % false: boxed links; true: colored links
linkcolor=red,          % color of internal links (change box color with linkbordercolor)
citecolor=blue,      % color of links to bibliography
filecolor=magenta,  % color of file links
urlcolor=cyan          % color of external links
}

\usepackage{array}

\externaldocument{../../SRS/SRS}

\input{../../Comments.text}
\input{../../Common.text}

\begin{document}

\title{Module Interface Specification for \progname{}}

\author{\authname}

\date{\today}

\maketitle

\pagenumbering{roman}

\section{Revision History}

\begin{tabularx}{\textwidth}{p{3cm}p{2cm}X}
\toprule {\bf Date} & {\bf Version} & {\bf Notes}\\
\midrule
Date 1 & 1.0 & Notes\\
Date 2 & 1.1 & Notes\\
\bottomrule
\end{tabularx}

~\newpage

\section{Symbols, Abbreviations and Acronyms}

See SRS Documentation at \wss{give url}

\wss{Also add any additional symbols, abbreviations or acronyms}

\newpage

\tableofcontents

\newpage

\pagenumbering{arabic}

\section{Introduction}

The following document details the Module Interface Specifications for
\wss{Fill in your project name and description}

Complementary documents include the System Requirement Specifications
and Module Guide.  The full documentation and implementation can be
found at \url{...}.  \wss{provide the url for your repo}

\section{Notation}

\wss{You should describe your notation.  You can use what is below as
  a starting point.}

The structure of the MIS for modules comes from \citet{HoffmanAndStrooper1995},
with the addition that template modules have been adapted from
\cite{GhezziEtAl2003}.  The mathematical notation comes from Chapter 3 of
\citet{HoffmanAndStrooper1995}.  For instance, the symbol := is used for a
multiple assignment statement and conditional rules follow the form $(c_1
\Rightarrow r_1 | c_2 \Rightarrow r_2 | ... | c_n \Rightarrow r_n )$.

The following table summarizes the primitive data types used by \progname. 

\begin{center}
\renewcommand{\arraystretch}{1.2}
\noindent 
\begin{tabular}{l l p{7.5cm}} 
\toprule 
\textbf{Data Type} & \textbf{Notation} & \textbf{Description}\\ 
\midrule
character & char & a single symbol or digit\\
integer & $\mathbb{Z}$ & a number without a fractional component in (-$\infty$, $\infty$) \\
natural number & $\mathbb{N}$ & a number without a fractional component in [1, $\infty$) \\
real & $\mathbb{R}$ & any number in (-$\infty$, $\infty$)\\
\bottomrule
\end{tabular} 
\end{center}

\noindent
The specification of \progname \ uses some derived data types: sequences, strings, and
tuples. Sequences are lists filled with elements of the same data type. Strings
are sequences of characters. Tuples contain a list of values, potentially of
different types. In addition, \progname \ uses functions, which
are defined by the data types of their inputs and outputs. Local functions are
described by giving their type signature followed by their specification.

\subsection{Shared Interface Types} \label{sec:interface-types}

The following records are shared by the session, records, dashboard, and
persistence interfaces. Each identifier type denotes an opaque identifier
for the named entity. An optional value is either present or absent.

\begin{description}
\item[StoreConfig] Configuration identifying a durable store.
\item[Actor] An authenticated user identity and permissions to act on patient
records or retrieve audit events.
\item[Patient] Patient identifier and identifying information.
\item[SpeechError] A record identifying a word occurrence in the practice
item and one of the following error kinds.
\begin{description}
\item[Substitution] Expected phoneme, produced phoneme, and expected
phoneme position within the word.
\item[Omission] Expected phoneme and its position within the word.
\item[Addition] Extra produced phoneme and its insertion location before,
between, or after expected phonemes. A location is the number of expected
phonemes preceding the insertion in that word. Multiple additions at the
same location retain their recognized order.
\end{description}
\item[Attempt] Attempt identifier, practice item, target phonemes, support
level, original analysis, and optional SLP correction. The analysis contains
expected and recognized IPA sequences (recognized IPA may be absent),
an ordered sequence of SpeechError, confidence when available, and
classification. It also contains an aligned observation
for each target occurrence, with an expected phoneme, optional produced
phoneme, and status (accepted, incorrect, or unclassifiable). The attempt
classification is correct, incorrect, or unclassifiable. An SLP correction
contains the SLP identity, optional correct/incorrect override, and any
phoneme-level corrections to the observations or error records.
\item[Session] Session identifier, patient identifier, SLP identity, session
date, session mode (testing or treatment), targets, practice items, SLP inputs, status (active or completed),
and a sequence of completed attempts. Only a completed snapshot is accepted
by \texttt{saveSession}.
\item[SessionPair] Two distinct session identifiers ordered from baseline
to follow-up.
\item[Progress] For a patient and target phoneme, dated session points
containing session identifiers and eligible, correct, and incorrect target
observations, with cue information where available. Includes an optional
comparison and optional improvement or persistent-error indicators.
\item[AuditEvent] Action, information type, timestamp, actor identity, and
retained patient identity reference.
\item[HistoryExport, AuditCopy] Copies of authorized session history and
audit events, respectively. The export formats are implementation decisions.
\end{description}

\section{Module Decomposition}

The following table shows the module hierarchy from the Module Guide.

\begin{table}[h!]
\centering
\begin{tabular}{p{0.3\textwidth} p{0.6\textwidth}}
\toprule
\textbf{Level 1} & \textbf{Level 2}\\
\midrule
{Hardware-Hiding Module} & ~ \\
\midrule
\multirow{6}{0.3\textwidth}{Behaviour-Hiding Module} & Audio Capture\\
& Speech Analysis\\
& Session Management\\
& Practice Items \& Lesson Plans\\
& Patient Records \& Progress Tracking\\
& Clinician Dashboard / Results Visualization\\
\midrule
{Software Decision Module} & Persistence\\
\bottomrule
\end{tabular}
\caption{Module Hierarchy}
\label{TblMH}
\end{table}

\clearpage
\section{MIS of Session Management} \label{mis:session}

\subsection{Module}
SessionManager runs one clinician-supervised practice session for one
patient, from setup to completion. It covers FR-SM1--4 and FR-PAR5--7, and the
session-side parts of FR-AC3, FR-SA3, FR-DMR2--3, FR-DMR6--7, and FR-DMR11.

\subsection{Uses}
Audio Capture supplies a recording and its usability flags (FR-AC1 and FR-AC3). Speech Analysis supplies the analysis of a usable recording (FR-SA1 and FR-SA3). Practice Items \& Lesson Plans supplies items, their target phonemes, and image cues. Patient Records \& Progress Tracking provides patient information if authorized and stores the completed session. Clinician Dashboard / Results Visualization uses this module; this module does not use the Dashboard.

\subsection{Syntax}

\subsubsection{Exported Types}
\begin{description}
\item[SessionMode] Testing or Treatment. Testing gathers information about the
patient's speech with no cues, hints, or modelling. Treatment allows support
from the SLP.
\item[SupportLevel] independent, semantic cue, or repetition/model. The support
a child received for one attempt (FR-SM3 and FR-DMR2.1).
\item[SessionSetup] Patient identifier, SessionMode, target phonemes, and
practice items.
\item[Prompt] A practice item and its image cue, presented for an attempt
(FR-PAR5).
\item[AttemptOutcome] Either \emph{flagged}, with the usability reasons from
Audio Capture, or \emph{analysed}, with the resulting Attempt.
\end{description}
Session, Attempt, Actor, and SLP correction are the shared types in
Section~\ref{sec:interface-types}. The Attempt records its SupportLevel.

\subsubsection{Exported Constants}
None.

\subsubsection{Exported Access Programs}
\noindent\begin{tabularx}{\textwidth}{@{}l X X X@{}}
\toprule
\textbf{Name} & \textbf{In} & \textbf{Out} & \textbf{Exceptions}\\
\midrule
\texttt{startSession} & Actor, SessionSetup & Session &
Denied, Missing, Invalid, Conflict, Unavailable\\
\texttt{startAttempt} & practice item & Prompt & Inactive, Invalid, Busy\\
\texttt{recordAttempt} & -- & AttemptOutcome & Inactive, NoAttempt,
Unavailable\\
\texttt{setSupport} & SupportLevel & -- & Inactive, NoAttempt, Invalid\\
\texttt{repeatAttempt} & -- & Prompt & Inactive, NoAttempt\\
\texttt{skipAttempt} & -- & -- & Inactive, NoAttempt\\
\texttt{overrideAttempt} & attempt identifier, SLP correction & -- &
Inactive, Missing, Invalid\\
\texttt{getSnapshot} & -- & Session & Inactive\\
\texttt{endSession} & -- & -- & Inactive, Denied, Invalid, Conflict,
Unavailable\\
\bottomrule
\end{tabularx}

\subsection{Semantics}

\subsubsection{State Variables}
\begin{itemize}
\item \texttt{session}: the active Session record, or absent. Initially absent.
Its attempts are the retained attempts so far, in order.
\item \texttt{actor}: the Actor who started the session.
\item \texttt{current}: the open attempt, or absent. It holds the practice
item, its SupportLevel (initially independent), and a state: prompted,
flagged, or analysed. An analysed current attempt is also the last element of
\texttt{session}'s attempts.
\item \texttt{nextId}: counter used to give attempts identifiers that are
unique within the session.
\end{itemize}
The session's SessionMode is part of the active Session record.

\subsubsection{Environment Variables}
A clock supplying the session date. The child-facing display, which shows the
Prompt.

\subsubsection{Assumptions}
The caller supplies an authenticated actor. One session is active at a time.
Each attempt is a single spoken target word (SRS A1), and the SLP supervises
the session (SRS A7). Audio Capture and Speech Analysis return the values
described in Section~\ref{sec:interface-types}. Recorded audio is used only
for its attempt (FR-AC2).

\subsubsection{Access Routine Semantics}
Every routine except \texttt{startSession} raises \texttt{Inactive} when no
session is active. Routines that act on the current attempt raise
\texttt{NoAttempt} when \texttt{current} is absent.

\begin{description}
\item[] \texttt{startSession(actor, setup)}\par
Raise \texttt{Conflict} if a session is already active. Confirm access to the
setup's patient identifier through Patient Records \& Progress Tracking;
\texttt{Denied}, \texttt{Missing}, and \texttt{Unavailable} propagate. Raise
\texttt{Invalid} when the target set or item set is empty, when a target is
outside the predetermined substitution list (FR-SM1 and FR-SA1), or when an
item does not exercise at least one selected target. On success, create an
active Session with a new identifier, the patient, the actor's SLP identity,
the setup's SessionMode, the clock date, the targets, the items, no SLP inputs,
and no attempts. Record the actor, clear \texttt{current}, and return the
Session.
\item[] \texttt{startAttempt(item)}\par
Raise \texttt{Invalid} if the item is not in the session's practice items.
Raise \texttt{Busy} if \texttt{current} is prompted or flagged and has not been
repeated or skipped. If \texttt{current} is analysed, it is final and stays in
the session. Set \texttt{current} to a prompted attempt for the item with
independent support, and return the Prompt supplied by Practice Items
(FR-PAR5).
\item[] \texttt{recordAttempt()}\par
Raise \texttt{Unavailable} if Audio Capture or Speech Analysis fails; the state
is unchanged. Obtain a recording from Audio Capture. If the recording is
flagged, set \texttt{current} to flagged, add nothing to the session, and
return flagged (FR-AC3). Otherwise obtain the analysis from Speech Analysis
for the item and its targets. Build an Attempt with a new identifier, the item,
its targets, the SupportLevel of \texttt{current}, and the original analysis,
which may classify it as unclassifiable (FR-SA3). Append it to the session's
attempts, set \texttt{current} to analysed, and return analysed. The Dashboard
sees the new attempt through \texttt{getSnapshot}.
\item[] \texttt{setSupport(level)}\par
Raise \texttt{Invalid} if the level is not in \texttt{allowedSupport(mode)},
where mode is the session's SessionMode. Set the SupportLevel of
\texttt{current}, and of its Attempt when analysed (FR-SM3 and FR-DMR2). The
SLP may call this before or after \texttt{recordAttempt}.
\item[] \texttt{repeatAttempt()}\par
Discard \texttt{current}. If it was analysed, remove its Attempt from the
session so that only the new attempt is retained. Start a new prompted attempt
for the same item with independent support and return its Prompt (FR-SM2).
\item[] \texttt{skipAttempt()}\par
Discard \texttt{current} as in \texttt{repeatAttempt}. Retain no Attempt for it
and set \texttt{current} to absent.
\item[] \texttt{overrideAttempt(attemptId, correction)}\par
Raise \texttt{Missing} if the identifier is not among the session's attempts.
Raise \texttt{Invalid} if the correction does not name the actor's SLP identity
or names a phoneme-level change for an observation that does not exist.
Attach the correction to the Attempt. The original analysis is not changed
(FR-SM4 and FR-DMR3). A new correction for the same attempt replaces the
previous one.
\item[] \texttt{getSnapshot()}\par
Return a copy of the Session with its current attempts. The Session's status is
active.
\item[] \texttt{endSession()}\par
Discard \texttt{current} if it is prompted or flagged. Form the completed
snapshot: the session with status completed. Call
\texttt{saveSession(actor, snapshot)} of Patient Records \& Progress Tracking.
On success, set \texttt{session} and \texttt{current} to absent. On any
exception (\texttt{Denied}, \texttt{Invalid}, \texttt{Conflict}, or
\texttt{Unavailable}), propagate it and leave the session active and its
attempts unchanged, so that calling \texttt{endSession} again submits an
identical snapshot. A session with zero attempts may be ended and saved.
\end{description}

\subsubsection{Local Functions}
\texttt{allowedSupport(mode)} returns \{independent\} for Testing and
\{independent, semantic cue, repetition/model\} for Treatment.

\noindent The effective classification of an attempt is defined by
\texttt{effectiveResult} in the Dashboard (Section~\ref{mis:dashboard}) and by
\texttt{summarize} in Patient Records \& Progress Tracking
(Section~\ref{mis:records}). This module does not compute it. The Dashboard
treats a correct attempt as cued when its SupportLevel is semantic cue or
repetition/model.

\clearpage
\section{MIS of Patient Records \& Progress Tracking} \label{mis:records}

\subsection{Module}
PatientRecords covers FR-DMR1--11 and FR-HA9, with Session Management supplying
session and attempt results.

\subsection{Uses}
Persistence supplies durable reads and atomic writes for patient records, sessions, and audit events.

\subsection{Syntax}

\subsubsection{Exported Types}
See the shared interface types in Section~\ref{sec:interface-types}.

\subsubsection{Exported Constants}
None.

\subsubsection{Exported Access Programs}
In the table, every operation except \texttt{initialize} may raise
\texttt{Unavailable} when the service has yet to initialize or persistence
fails, and \texttt{Denied} if authorization fails.

\noindent\begin{tabularx}{\textwidth}{@{}l X X l@{}}
\toprule
\textbf{Name} & \textbf{In} & \textbf{Out} & \textbf{Other Exceptions}\\
\midrule
\texttt{initialize} & StoreConfig & -- & Unavailable\\
\texttt{createPatient} & Actor, Patient & -- & Conflict\\
\texttt{updatePatient} & Actor, Patient & -- & Missing\\
\texttt{getPatient} & Actor, patient identifier & Patient & Missing\\
\texttt{deletePatient} & Actor, patient identifier, Boolean confirmation & -- &
Missing, Unconfirmed\\
\texttt{saveSession} & Actor, completed Session & -- & Missing, Invalid, Conflict\\
\texttt{getHistory} & Actor, patient identifier & Sequence of Session & Missing\\
\texttt{exportHistory} & Actor, patient identifier & HistoryExport & Missing\\
\texttt{getProgress} & Actor, patient identifier, target phoneme,
optional SessionPair & Progress & Missing\\
\texttt{getAudit} & Actor, retained patient identity reference &
Sequence of AuditEvent & --\\
\texttt{copyAudit} & Actor, retained patient identity reference &
AuditCopy & --\\
\bottomrule
\end{tabularx}

\subsection{Semantics}
\subsubsection{State Variables}
An initialization flag, initially false, and a reference to the configured
Persistence store after initialization. Patient data and audit events are
maintained by Persistence.

\subsubsection{Environment Variables}
A clock supplying timestamps for audit events.

\subsubsection{Assumptions}
The caller supplies an authenticated actor. Access to a patient's records and
access to the audit trail are independently authorized.

\subsubsection{Access Routine Semantics}
For all patient operations, evaluate authorization before returning patient
data or writing it. Raise \texttt{Missing} when an authorized request refers
to a patient absent from active records. Record successful create, view,
modify, export, and delete actions with their information type, actor,
timestamp, and patient reference. If the audit write fails, raise
\texttt{Unavailable} before exposing the output or committing a mutation.

\begin{description}
\item[] \texttt{initialize(store)}\par
Set the store and initialization flag on
success. On failure, raise \texttt{Unavailable}.
\item[] \texttt{createPatient(actor, patient)}\par
Insert the patient and creation
audit event. Raise \texttt{Conflict} for a duplicate identifier.
\item[] \texttt{updatePatient(actor, patient)}\par
Replace identifying information
for the given identifier and append a modification event.
\item[] \texttt{getPatient(actor, patientId)}\par
Append a view event and return the
patient record.
\item[] \texttt{deletePatient(actor, patientId, confirmed)}\par
Raise
\texttt{Unconfirmed} if the SLP has not confirmed deletion for that patient.
Otherwise remove the patient, sessions, and attempts from active records in
one operation. Preserve prior audit events and append a deletion event
under the retained patient identity reference.
\item[] \texttt{saveSession(actor, session)}\par
Validate completion, unique attempt
identifiers, and the patient's association. Raise \texttt{Invalid} for a
snapshot that fails validation. Commit the session, attempts, and modification
event together. An identical retry leaves them unchanged. Raise
\texttt{Conflict} when the session identifier already denotes a different
snapshot. The original analysis remains intact when an override is present.
\item[] \texttt{getHistory(actor, patientId)}\par
Append a view event and return
that patient's completed sessions ordered by date, with ties ordered by
session identifier. A patient with no sessions has empty history.
\item[] \texttt{exportHistory(actor, patientId)}\par
Append an export event and
return a copy of the same ordered history.
\item[] \texttt{getProgress(actor, patientId, target, pair)}\par
Compute progress
from \texttt{summarize(history, target)} for this patient's history. If a
pair is supplied, compare its session points. Raise \texttt{Missing} if
either session is outside this patient's history. Return available session
points with an absent comparison if either selected session has no eligible
observations for the target. If no pair is supplied, return the history
without a between-session comparison. Append a view event before returning.
\item[] \texttt{getAudit(actor, patientRef)} and \texttt{copyAudit(actor, patientRef)}\par
Return matching events in append order or a copy of them. Include retained
events for deleted patients when the actor has audit access.
\end{description}

\subsubsection{Local Functions}
\texttt{summarize(history, target)} reports per-session target counts and
support levels. Use phoneme-level corrections where present. Count accepted
productions as correct and substitutions or omissions as incorrect. Exclude
additions and unresolved unclassifiable occurrences. Attempt-level overrides
leave phoneme counts unchanged.

The correct rate is correct divided by eligible occurrences. Compare sessions
as follow-up minus baseline. Use the two latest eligible sessions for trends:
a higher rate indicates improvement, and errors in both indicate persistence.
Fewer than two eligible sessions yields no comparison or trend.

\clearpage
\section{MIS of Clinician Dashboard / Results Visualization} \label{mis:dashboard}

\subsection{Module}
Template ClinicianDashboard. One instance presents results for a clinician
(FR-RDV1--9, FR-DMR8--10, and FR-HA7).

\subsection{Uses}
Session Management supplies the selected session
snapshot and effective clinician overrides. Patient Records \& Progress
Tracking (Section~\ref{mis:records}) supplies authorized longitudinal results.

\subsection{Syntax}

\subsubsection{Exported Types}
\texttt{Dashboard} denotes a clinician view instance. Its result panels
present attempt detail, a session summary and heat map, and per-target
progress. Identifiers and snapshots use the types in
Section~\ref{sec:interface-types}.

\subsubsection{Exported Constants}
None.

\subsubsection{Exported Access Programs}
\noindent\begin{tabularx}{\textwidth}{@{}l X X l@{}}
\toprule
\textbf{Name} & \textbf{In} & \textbf{Out} & \textbf{Exceptions}\\
\midrule
\texttt{newDashboard} & Actor & Dashboard & Denied\\
\texttt{showAttempt} & Dashboard, Session, attempt identifier & -- &
Denied, Missing\\
\texttt{showSession} & Dashboard, Session & -- & Denied\\
\texttt{showProgress} & Dashboard, patient identifier, target phoneme,
optional SessionPair & -- & Denied, Missing, Unavailable\\
\bottomrule
\end{tabularx}

\subsection{Semantics}

\subsubsection{State Variables}
The dashboard's actor, optional selected patient identifier, and displayed
panels. Initially, the selection and panels are empty. All visible patient
panels belong to the selected patient.

\subsubsection{Environment Variables}
The clinician display, updated when a result panel changes.

\subsubsection{Assumptions}
Session Management supplies authorized snapshots and a refreshed snapshot
after an attempt completes, an override changes, or a session ends.
Phoneme closeness uses the SRS phonological feature vectors. Recognition
confidence is displayed separately.

\subsubsection{Access Routine Semantics}
When a request selects another patient, clear the previous patient's panels.
Raise \texttt{Denied} and clear patient content if the actor lacks access.
Preface results with an acknowledgement that the information is assistive
and may be incorrect, and that interpretation is at the SLP's discretion.
Use neutral wording on clinician result panels.

\begin{description}
\item[] \texttt{newDashboard(actor)}\par
Return an empty dashboard for an authorized
clinician. Raise \texttt{Denied} for other actors.
\item[] \texttt{showAttempt(view, session, attemptId)}\par
Display expected and recognized IPA side by side. For each error, show the
expected-to-produced substitution, the omitted expected phoneme with a gap
in the recognized IPA, or the added produced phoneme with a gap in the
expected IPA. Show the beginning, middle, or end word position for a
substitution or omission and the insertion location for an addition. Raise
\texttt{Missing} if the attempt is absent from the supplied session. For an
unclassifiable attempt, show its status and any recognized IPA that is
available. When the SLP overrides a classification, show the effective
classification, identify it as the SLP's decision, and keep the original
analysis visible.
\item[] \texttt{showSession(view, session)}\par
Recompute the view from the supplied
snapshot. Count correct, incorrect, and unresolved unclassifiable attempts
using \texttt{effectiveResult}. Report substitution, omission, and addition
counts separately, using corrected error records when supplied. Split
correct attempts by independent and
cued support. Show zero counts and an empty heat map for zero attempts, and
leave rates with a zero denominator unavailable. Update the open attempt
panel from the same snapshot if applicable. Apply \texttt{closeness} to
eligible expected-to-produced pairs. A classification-only override does
not change the pair, and unresolved unclassifiable results without an
eligible pair have no heat-map value. Show omissions and additions as
unscored gaps in the IPA comparison.
\item[] \texttt{showProgress(view, patientId, target, pair)}\par
Retrieve and display
authorized progress for the given patient and target in date order.
Display the selected pair's comparison and the trend indicators when
available. If the target has fewer than two eligible session points or
the selected pair lacks eligible observations, show the available points
and indicate insufficient history. On
\texttt{Missing} or \texttt{Unavailable}, clear the requested panel and
report the exception.
\end{description}

\subsubsection{Local Functions}
\texttt{effectiveResult(attempt)} returns the SLP classification override
when present and the original classification otherwise. A correct attempt
with any additional cue or repetition is cued. A correct attempt recorded as
independent is independent.

\noindent\texttt{closeness(expected, produced)} returns
$1 - \frac{\sum_{i=1}^{n}|A_i-B_i|}{2n}$ for the SRS feature vectors
$A$ and $B$, each containing $n>0$ features in $\{-1,0,1\}$.
Use a colour scale from 0 (least close) to 1 (closest). Use the corrected
produced phoneme when the SLP supplies one. A pair with a missing feature
vector or no classifiable or corrected produced phoneme has no heat-map
value.

\bibliographystyle {plainnat}
\bibliography {../../../refs/References}

\newpage

\newpage{}

\section{Appendix --- Generative AI Use Disclosure}

\wss{ChatGPT (version 5) was used to brainstorm ideas for the template for this disclose section.}

\subsection{AI Tools Used}

\wss{Give the name of the AI tool(s) used and the version.}

\subsection{How and Where Used}

\wss{Explain what the tool(s) were used for.  Examples include: brainstorming,
explaining a technical concept, reviewing the document, generating or modifying
text, generating or modifying diagrams, identifying errors or omissions or
inconsistencies.}

\wss{For each of the uses, identify which parts of the document were affected.}

\wss{You do not need to report every interaction, but you should disclose the 
uses that materially contributed to your work}

\subsection{Verification of AI-Generated Content}

\wss{\begin{itemize}
    \item Describe how the team checked AI-generated or AI-assisted material
    before including it in the document.
    \item In particular, verify factual claims, technical assertions,
    requirements, calculations, references, and other information for which an
    AI system may produce incorrect or fabricated results.
    \item \textbf{The team is responsible for the accuracy, completeness,
consistency, and appropriateness of the submitted document, regardless of
whether material was generated by a human or an AI system.}
\end{itemize}}

\newpage{}

\section{Appendix --- Reflection}

\wss{Not required for CAS 741 projects}

The information in this section will be used to evaluate the team members on the
graduate attribute of Problem Analysis and Design.

\input{../../Reflection.text}

\begin{enumerate}
  \item What went well while writing this deliverable? 
  \item What pain points did you experience during this deliverable, and how
    did you resolve them?
  \item Which of your design decisions stemmed from speaking to your client(s)
  or a proxy (e.g. your peers, stakeholders, potential users)? For those that
  were not, why, and where did they come from?
  \item While creating the design doc, what parts of your other documents (e.g.
  requirements, hazard analysis, etc), it any, needed to be changed, and why?
  \item What are the limitations of your solution?  Put another way, given
  unlimited resources, what could you do to make the project better? (LO\_ProbSolutions)
  \item Give a brief overview of other design solutions you considered.  What
  are the benefits and tradeoffs of those other designs compared with the chosen
  design?  From all the potential options, why did you select the documented design?
  (LO\_Explores)
  \item (After you have implemented another team's module, which means this
  isn't filled in until after the original deadline). What did you learn by
  implementing another team's module? Were all the details you needed in the
  documentation, or did you need to make assumptions, or ask the other team
  questions?  If your team also had another team implement one of your modules,
  what was this experience like?  Are there things in your documentation you
  could have changed to make the process go more smoothly for when an
  ``outsider'' completes some of the implementation?

  Your reflection should be about concepts learned in your SE education, such as modular decomposition, cohesion, coupling, information hiding, etc.  You should critique the quality and utility of the specification you received.

  You should create a PR for your implemented module.  A link to the PR should be part of your reflection.

\end{enumerate}

\end{document}
